\documentclass[12pt,a4paper]{article}
\usepackage[utf8]{inputenc}
\usepackage[spanish]{babel}
\usepackage{amsmath}
\usepackage{amsfonts}
\usepackage{amssymb}
\usepackage{physics}
\usepackage{listings}
\usepackage{xcolor}
\usepackage{graphicx}
\usepackage{hyperref}
\usepackage{geometry}
\usepackage{siunitx}
\usepackage{booktabs}

\geometry{margin=2.5cm}

\lstset{
    language=C,
    basicstyle=\ttfamily\small,
    keywordstyle=\color{blue},
    commentstyle=\color{green!60!black},
    stringstyle=\color{red},
    numberstyle=\tiny\color{gray},
    breaklines=true,
    frame=single,
    numbers=left
}

\title{\textbf{Análisis Académico de una Simulación de Gas Unidimensional}\\
       \large{Aspectos Físicos y Computacionales}}
\author{Análisis del código OpenMP C}
\date{\today}

\begin{document}

\maketitle

\tableofcontents
\newpage

\section{Introducción}

Este documento presenta un análisis académico exhaustivo de una simulación computacional de partículas de gas en una dimensión espacial. La simulación emplea OpenMP para paralelización y utiliza métodos de Monte Carlo para modelar el comportamiento estocástico del sistema. Analizaremos tanto los aspectos físicos fundamentales como los aspectos computacionales de la implementación.

\section{Marco Teórico Físico}

\subsection{Modelo Físico Fundamental}

La simulación implementa un modelo de gas unidimensional basado en los siguientes principios:

\subsubsection{Ecuación de Movimiento Clásica}

El movimiento de cada partícula sigue la ecuación newtoniana básica:

\begin{equation}
x_{n+1} = x_n + \frac{p_n \cdot \Delta t}{m}
\end{equation}

donde:
\begin{itemize}
    \item $x_n$ es la posición de la partícula en el paso temporal $n$
    \item $p_n$ es el momento lineal de la partícula
    \item $m$ es la masa de la partícula (\SI{6.646473667973e-27}{\kilogram})
    \item $\Delta t$ es el paso temporal (\SI{4.430982591982e-7}{\second})
\end{itemize}

Esta ecuación representa la integración temporal de primer orden de la velocidad clásica $v = p/m$.

\subsubsection{Condiciones de Frontera Periódicas con Reflexión Estocástica}

El sistema implementa condiciones de frontera especiales. Cuando una partícula alcanza los límites del dominio espacial ($|x| > 0.5$), se produce un evento de reflexión que incluye:

\begin{equation}
k = \text{trunc}(x + 0.5 \cdot \text{sign}(p))
\end{equation}

donde $k$ representa el número de cruces de frontera. Para cada cruce ($k \neq 0$), se aplica:

\begin{equation}
\Delta x = \sqrt{|k|} \cdot \xi_1 \cdot \cos(\xi_2) \cdot \sigma_L
\end{equation}

donde $\xi_1$ y $\xi_2$ son variables aleatorias gaussiana y uniforme respectivamente, y $\sigma_L = \SI{0.0001}{\meter}$ es la incertidumbre en la posición.

\subsubsection{Modificación Estocástica del Momento}

Durante cada reflexión, el momento se modifica según:

\begin{equation}
\Delta E = \alpha \cdot (|p| - p_{\min}) \cdot (p_{\max} - |p|)
\end{equation}

\begin{equation}
p_{\text{new}} = \sqrt{\max(0, p^2 + \Delta E \cdot (\text{random} - 0.5))}
\end{equation}

donde:
\begin{itemize}
    \item $\alpha$ es un parámetro de acoplamiento (inicializado en 0)
    \item $p_{\min} = \SI{2.0e-26}{\kilogram\meter\per\second}$
    \item $p_{\max} = \SI{3.0e-23}{\kilogram\meter\per\second}$
\end{itemize}

\subsection{Distribución de Equilibrio}

\subsubsection{Distribución Espacial}

El sistema espera una distribución espacial uniforme en el estado de equilibrio:

\begin{equation}
\rho_x(x) = \frac{N_{\text{part}}}{L}
\end{equation}

donde $L = 1$ es la longitud del dominio espacial y $N_{\text{part}} = 2^{21} = 2,097,152$ es el número total de partículas.

\subsubsection{Distribución de Maxwell-Boltzmann}

La distribución esperada de momentos sigue una gaussiana:

\begin{equation}
\rho_p(p) = \frac{N_{\text{part}}}{\sigma_p \sqrt{2\pi}} \exp\left(-\frac{(p - \mu_p)^2}{2\sigma_p^2}\right)
\end{equation}

donde $\sigma_p = \SI{5.24684e-24}{\kilogram\meter\per\second}$ y $\mu_p = 0$ (distribución centrada).

\section{Implementación Computacional}

\subsection{Arquitectura de Datos}

\subsubsection{Tipos de Datos Primitivos}

La simulación utiliza tipos de datos específicos para optimización:

\begin{lstlisting}[caption=Tipos de datos principales]
double *x;    // Posiciones: IEEE 754 double precision (64 bits)
double *p;    // Momentos: IEEE 754 double precision (64 bits)
int *h;       // Histograma espacial: enteros con signo de 32 bits
int *g;       // Histograma de momentos: enteros con signo de 32 bits
uint32_t;     // Semillas PRNG: enteros sin signo de 32 bits
\end{lstlisting}

\subsubsection{Estructura de Memoria}

El programa aloca memoria de forma dinámica para arrays unidimensionales:

\begin{equation}
\begin{aligned}
\text{Memoria}_{\text{posiciones}} &= N_{\text{part}} \times 8 \text{ bytes} \\
\text{Memoria}_{\text{momentos}} &= N_{\text{part}} \times 8 \text{ bytes} \\
\text{Memoria}_{\text{histogramas}} &= (4 \times \text{BINS} + 4) \times 4 \text{ bytes}
\end{aligned}
\end{equation}

Para $N_{\text{part}} = 2^{21}$ y $\text{BINS} = 500$:
- Memoria total para partículas: $\sim 33.6$ MB
- Memoria para histogramas: $\sim 8$ KB

\subsection{Algoritmos de Generación de Números Aleatorios}

\subsubsection{Generador XORShift}

La simulación implementa el algoritmo XORShift para generación de números pseudoaleatorios:

\begin{lstlisting}[caption=Implementación XORShift]
inline double d_xorshift(uint32_t *state) {
    uint32_t x = *state;
    x ^= x << 13;    // Operación XOR con desplazamiento izquierdo
    x ^= x >> 17;    // Operación XOR con desplazamiento derecho
    x ^= x << 5;     // Operación XOR con desplazamiento izquierdo
    *state = x;
    return (double)x / (double)UINT32_MAX;
}
\end{lstlisting}

Este algoritmo tiene período $2^{32} - 1$ y es computacionalmente eficiente para simulaciones paralelas.

\subsubsection{Transformación de Box-Muller}

Para generar variables aleatorias gaussianas, se utiliza la transformación de Box-Muller:

\begin{equation}
\begin{aligned}
\xi_1 &= \sqrt{-2\ln(U_1 + 10^{-35})} \\
\xi_2 &= 2\pi U_2 \\
Z_1 &= \xi_1 \cos(\xi_2) \\
Z_2 &= \xi_1 \sin(\xi_2)
\end{aligned}
\end{equation}

donde $U_1, U_2$ son variables uniformes en $[0,1]$ y $Z_1, Z_2$ son gaussianas independientes con $\mu = 0, \sigma = 1$.

\subsection{Paralelización con OpenMP}

\subsubsection{Estrategias de Paralelización}

El código emplea varias directivas OpenMP para optimizar el rendimiento:

\begin{lstlisting}[caption=Paralelización de la inicialización]
#pragma omp parallel for reduction(+:DpE[:2*BINS]) schedule(static)
for (int i = 0; i < BINS << 1; i++) {
    // Cálculo de distribución esperada de momentos
}
\end{lstlisting}

\begin{lstlisting}[caption=Paralelización del bucle principal]
#pragma omp parallel shared(x, p)
{
    uint32_t seed = (uint32_t)(time(NULL) + omp_get_thread_num());
    #pragma omp for private(k, signop) schedule(dynamic)
    for (int i = 0; i < N_PART; ++i) {
        // Evolución temporal de cada partícula
    }
}
\end{lstlisting}

\subsubsection{Análisis de Complejidad Computacional}

\begin{itemize}
    \item \textbf{Complejidad temporal}: $O(N_{\text{part}} \times N_{\text{steps}})$ por iteración
    \item \textbf{Complejidad espacial}: $O(N_{\text{part}} + \text{BINS})$
    \item \textbf{Escalabilidad}: Lineal con el número de threads hasta límite de memoria compartida
\end{itemize}

\section{Análisis Estadístico}

\subsection{Pruebas de Bondad de Ajuste}

El programa implementa múltiples pruebas chi-cuadrado para validar las distribuciones:

\subsubsection{Chi-cuadrado para Distribución Espacial}

\begin{equation}
\chi^2_x = \frac{1}{2\text{BINS} - 4} \sum_{i=4}^{2\text{BINS}} \frac{(h_i - D_{x,i})^2}{D_{x,i}}
\end{equation}

donde $h_i$ es el histograma observado y $D_{x,i}$ es la distribución teórica esperada.

\subsubsection{Chi-cuadrado para Distribución de Momentos}

\begin{equation}
\chi^2_p = \frac{1}{2(\text{BINS} - \text{BORDES})} \sum_{i=\text{BORDES}}^{2\text{BINS}-\text{BORDES}} \frac{(g_i - D_{p,i})^2}{D_{p,i}}
\end{equation}

\subsubsection{Pruebas de Simetría}

Para verificar la simetría espacial:

\begin{equation}
\chi^2_{I,x} = \frac{1}{2(\text{BINS} - 2)} \sum_{i=4}^{\text{BINS}+1} \frac{(h_i - h_{2\text{BINS}+3-i})^2}{D_{x,i}}
\end{equation}

\section{Conservación de Energía}

\subsection{Cálculo de Energía Total}

La energía cinética total del sistema se calcula como:

\begin{equation}
E_{\text{total}} = \sum_{i=1}^{N_{\text{part}}} \frac{p_i^2}{2m}
\end{equation}

Esta cantidad se monitorea durante toda la simulación para verificar la conservación aproximada de energía.

\subsection{Implementación Paralela del Cálculo Energético}

\begin{lstlisting}[caption=Cálculo paralelo de energía]
#pragma omp parallel for reduction(+:sumEnergy) schedule(static)
for (int i = 0; i < N_PART; i++) {
    sumEnergy += p[i] * p[i];
}
return sumEnergy / (2 * M);
\end{lstlisting}

La directiva \texttt{reduction} asegura que la suma se realice de forma thread-safe.

\section{Gestión de Archivos y Persistencia}

\subsection{Formato de Archivos de Datos}

Los archivos de estado del sistema utilizan formato binario:

\begin{lstlisting}[caption=Estructura de archivos binarios]
struct {
    unsigned int evolution;     // 4 bytes: número de pasos
    double x[N_PART];          // 8*N_PART bytes: posiciones
    double p[N_PART];          // 8*N_PART bytes: momentos
};
\end{lstlisting}

\subsection{Archivos de Histogramas}

Los histogramas se guardan en formato ASCII con la siguiente estructura:

\begin{equation}
\text{Formato: } x_i \quad h_i \quad p_i \quad g_i \quad \text{(estadísticas chi-cuadrado)}
\end{equation}

\section{Optimizaciones y Consideraciones Numéricas}

\subsection{Precisión Numérica}

\begin{itemize}
    \item \textbf{Punto flotante}: IEEE 754 double precision (15-17 dígitos decimales)
    \item \textbf{Estabilidad temporal}: $\Delta t$ elegido para mantener $\frac{p \Delta t}{m} \ll 1$
    \item \textbf{Condición CFL}: Implícitamente satisfecha por la elección de parámetros
\end{itemize}

\subsection{Técnicas de Optimización}

\begin{enumerate}
    \item \textbf{Operaciones bit-shift}: \texttt{BINS << 1} en lugar de \texttt{2*BINS}
    \item \textbf{Scheduling dinámico}: Para balancear carga entre threads
    \item \textbf{Memory alignment}: Arrays alocados dinámicamente para evitar cache misses
    \item \textbf{Loop unrolling}: Implícito en compiladores optimizadores
\end{enumerate}

\section{Validación Física}

\subsection{Teorema de Equipartición}

En equilibrio térmico, el sistema debe satisfacer:

\begin{equation}
\langle E_{\text{cinética}} \rangle = \frac{1}{2} k_B T
\end{equation}

donde $k_B$ es la constante de Boltzmann y $T$ la temperatura efectiva del sistema.

\subsection{Criterios de Convergencia}

La simulación considera convergencia cuando:

\begin{equation}
\chi^2_x < \text{umbral} \quad \text{y} \quad \chi^2_p < \text{umbral}
\end{equation}

Típicamente, umbrales alrededor de 1.0-2.0 indican buen ajuste a las distribuciones teóricas.

\section{Conclusiones}

Esta simulación representa una implementación robusta de dinámica molecular unidimensional con las siguientes características destacadas:

\begin{enumerate}
    \item \textbf{Física}: Implementa correctamente ecuaciones newtonianas con condiciones de frontera estocásticas
    \item \textbf{Computación}: Utiliza paralelización eficiente y algoritmos numericos estables
    \item \textbf{Estadística}: Incluye validación rigurosa mediante pruebas chi-cuadrado
    \item \textbf{Escalabilidad}: Arquitectura preparada para sistemas de gran escala
\end{enumerate}

El código demuestra una excelente integración entre principios físicos fundamentales y técnicas computacionales avanzadas, proporcionando una herramienta valiosa para el estudio de sistemas de muchas partículas en una dimensión.

\section{Referencias Técnicas}

\begin{enumerate}
    \item OpenMP Architecture Review Board. \textit{OpenMP Application Programming Interface}, versión 4.5
    \item Marsaglia, G. \textit{Xorshift RNGs}. Journal of Statistical Software, 2003
    \item Box, G.E.P.; Muller, M.E. \textit{A Note on the Generation of Random Normal Deviates}. Annals of Mathematical Statistics, 1958
    \item Allen, M.P.; Tildesley, D.J. \textit{Computer Simulation of Liquids}. Oxford University Press, 1987
\end{enumerate}

\end{document} 